% Source pp. 201-207 (Lectures19-25.pdf).
\section{Some Properties of RG Flow}

In the previous Section we have discussed the basic properties of the Gaussian
fixed point $\mathcal A_{*G}$ at different dimensions $d$. At special values of
$d$ the fixed point theory has marginal operators. Thus, at $d=4$ the field
$\Phi_4\sim\varphi^4$ is marginal, and as more careful calculation shows, it is
\emph{marginally irrelevant}. This is expressed by the RG flow equation
\eqref{eq:WilsonBeta} (note that it is equivalent to the leading term in the beta
function \eqref{eq:beta1}).

\subsection*{RG flow and Poincar\'e-Dulac normal form}

It is interesting to see, in general, what is special about the marginal
deformations from the mathematical point of view. Discussing on the abstract
level, let us assume that we have generic critical fixed point
$\mathcal A_*\in\Sigma$. Let $y^\alpha$ be some local coordinates on $\Sigma$,
covering some neighborhood of $\mathcal A_*$ and chosen in such a way that the
fixed point corresponds to $y^\alpha=0$. One can think of $y^\alpha$ as the
collection of coupling constants $u_{2k}-u_{2k,*}$, $v_{2k}-v_{2k,*}$, $\dots$,
where $u$'s, $v$'s, etc, are the coefficients appearing in writing $\mathcal A$ in
terms of various monomials in $\varphi_0$ and its derivatives, and $u_*$'s, etc,
are their fixed-point values. Anyhow, the RG flow equation
\[
\frac d{dl}\,\mathcal A=B\{\mathcal A\}
\]
can be written in such coordinates as
\begin{equation}\label{eq:flowy}
\frac d{dl}\,y^\alpha=B^\alpha(\{y\})\,,
\end{equation}
where $B^\alpha$ are some functions of all $y$'s. As $y=0$ is assumed to be a fixed
point, these functions satisfy the condition
\[
B^\alpha(\{0\})=0\,.
\]
Let us assume that these functions can be expanded into Taylor series around the
fixed point,
\[
B^\alpha(\{y\})=K^\alpha_\beta\,y^\beta+C^\alpha_{\beta\gamma}\,y^\beta y^\gamma+\cdots\,,
\]
where $K$, $C$, etc, are some constants, and summation over repeated indices is
assumed. The coefficients $K^\alpha_\beta$ are the matrix elements of the operator
$K$, the linearization of the RG generator $B$ near the fixed point,
\[
B\{\mathcal A_*+\delta\mathcal A\}=K\,\delta\mathcal A+O(\delta\mathcal A^2)\,.
\]
By appropriate linear transformation of the coordinates $y$:
$y^\alpha\to M^\alpha_\beta y^\beta$ the matrix $K^\alpha_\beta$ can be brought to
diagonal form.\footnote{Non-diagonalizable forms contradict unitarity (positivity
of norms).} In these coordinates the flow equation \eqref{eq:flowy} takes the form
\begin{equation}\label{eq:flowdiag}
\frac d{dl}\,y^\alpha=\kappa_\alpha\,y^\alpha+c^\alpha_{\beta\gamma}\,y^\beta y^\gamma+\cdots\,,
\end{equation}
where $\kappa_\alpha$ are the eigenvalues of $K$, and no summation over $\alpha$ is
assumed in the first term in the r.h.s.; the coefficients $c$ are related to $C$
by a linear transformation, in an obvious way.

No further simplification of the linear term in the r.h.s.\ of \eqref{eq:flowdiag}
is possible. However one may ask if the quadratic and higher order terms in the
r.h.s.\ can be simplified by nonlinear coordinate transformation
\[
y^\alpha=y^\alpha(\{\tilde y\})=\tilde y^\alpha+a^\alpha_{\beta\gamma}\,\tilde y^\beta\tilde y^\gamma+\cdots\,.
\]
Under coordinate transformations the functions $B^\alpha$ transform as components
of a vector field on $\Sigma$. Simple calculation shows that in terms of $\tilde y$
the RG flow equations take the form
\[
\frac d{dl}\,\tilde y^\alpha=\kappa_\alpha\,\tilde y^\alpha+\tilde c^\alpha_{\beta\gamma}\,\tilde y^\beta\tilde y^\gamma+\cdots\,,
\qquad\text{where}\qquad
\tilde c^\alpha_{\beta\gamma}=c^\alpha_{\beta\gamma}+(\kappa_\alpha-\kappa_\beta-\kappa_\gamma)\,a^\alpha_{\beta\gamma}\,.
\]
One can completely eliminate the quadratic term by taking
\[
a^\alpha_{\beta\gamma}=-\frac{c^\alpha_{\beta\gamma}}{\kappa_\alpha-\kappa_\beta-\kappa_\gamma}\,;
\]
with this choice
\[
\frac d{dl}\,\tilde y^\alpha=\kappa_\alpha\,\tilde y^\alpha+O(\tilde y^3)\,.
\]
This is true provided the combinations of the eigenvalues appearing in the
denominators do not vanish. However if some triplet of the eigenvalues satisfies
the \emph{resonance condition}
\[
\kappa_\alpha-\kappa_\beta-\kappa_\gamma=0
\]
the ``resonance'' term
\[
c^\alpha_{\beta\gamma}\,y^\beta y^\gamma
\]
(no summation) cannot be eliminated.

Similar situation takes place in the higher orders. If components $B^\alpha$ of a
vector field $B$ can be written as power series
\begin{equation}\label{eq:Bseries}
B^\alpha(\{y\})=\kappa_\alpha\,y^\alpha+c^\alpha_{\beta_1\beta_2}\,y^{\beta_1}y^{\beta_2}+c^\alpha_{\beta_1\beta_2\beta_3}\,y^{\beta_1}y^{\beta_2}y^{\beta_3}+\cdots\,.
\end{equation}
in some local coordinates $y^\alpha$, then by means of appropriate analytic
coordinate transformation
\begin{equation}\label{eq:normaltr}
y^\alpha\to y^\alpha+a^\alpha_{\beta_1\beta_2}\,y^{\beta_1}y^{\beta_2}+a^\alpha_{\beta_1\beta_2\beta_3}\,y^{\beta_1}y^{\beta_2}y^{\beta_3}+\cdots
\end{equation}
one can eliminate all nonlinear terms in \eqref{eq:Bseries}, except for the
so-called resonance terms. A term
\begin{equation}\label{eq:resterm}
c^\alpha_{\beta_1\beta_2\cdots\beta_n}\,y^{\beta_1}y^{\beta_2}\cdots y^{\beta_n}
\end{equation}
in the function $B^\alpha$ is called the \emph{resonance term} if the following
\emph{resonance condition}
\begin{equation}\label{eq:rescond}
\kappa_\alpha=\kappa_{\beta_1}+\kappa_{\beta_2}+\dots+\kappa_{\beta_n}
\end{equation}
is satisfied. One can always adjust the coefficients $a$ in \eqref{eq:normaltr},
order by order in $y$, in such a way that in new coordinates
\begin{equation}\label{eq:normalform}
\frac d{dl}\,y^\alpha=\kappa_\alpha\,y^\alpha+\text{Resonance Terms}
\end{equation}
If there are no resonances, all nonlinear terms in \eqref{eq:Bseries} can be
eliminated. This statement is known as \emph{Poincar\'e-Dulac
theorem}.\footnote{See e.g.\ V.I.~Arnold, Geometrical Methods in the Theory of
Ordinary Differential Equations, Springer, 1988.} I will call ``canonical'' the
coordinates $y$ in which the flow equation takes the form \eqref{eq:normalform}.

Two remarks are in order. First, the above theorem \emph{does not} suggest that
the global structure of the RG flow is in any sense ``trivial'', even in
non-resonance cases. The series \eqref{eq:normaltr} which defines the
transformation converges at best in some finite neighborhood of the fixed point
$y=0$. Other interesting irregularities the global RG flow may possess (like
other fixed points) are located no closer than at the boundary of this domain of
convergence; in the new coordinates where $B^\alpha$'s are ``trivial'' they are
``pushed out to the infinity''. To get some idea consider the following

\begin{exercise}
Study the flow equation with just one coordinate,
\[
\frac d{dl}\,y=y-y^2
\]
(which has two fixed points $y=0$ and $y=1$). Find transformation
$y=f(\tilde y)=\tilde y+\cdots$ such that
\[
\frac d{dl}\,\tilde y=\tilde y\,.
\]
Study the properties of the map $y\to\tilde y$.
\end{exercise}

Second remark concerns very simple interpretation which the above resonance
conditions have in terms of ordinary dimensional analysis. In the neighborhood of
$\mathcal A_*$ where the coordinates $y$ are defined, the action can be written as
\[
\mathcal A=\mathcal A_*+y^\alpha\int d^dx\,\Phi_\alpha(x)+O(y^2)\,,
\]
where $\Phi_\alpha(x)$ are the fields which diagonalize the operator $D$
associated with the fixed point $\mathcal A_*$. The eigenvalues $D_\alpha$ (i.e.\
the dimensions of the fields $\Phi_\alpha$) are related to $\kappa_\alpha$ as
\[
D_\alpha+\kappa_\alpha=d\,.
\]
Therefore the numbers $\kappa_\alpha$ can be interpreted as the mass dimensions of
the ``coupling constants'' $y^\alpha$. The resonance condition \eqref{eq:rescond}
then states that the resonance term \eqref{eq:resterm} has ``right dimension''. In
this sense, the coordinate system in which the RG flow equations have the form
\eqref{eq:normalform} respects the straightforward dimensional analysis.

From this point of view, what distinguishes the marginal case is that the presence
of a marginal field (i.e.\ the field $\Phi_\alpha$ with $D_\alpha=d$, or
$\kappa_\alpha=0$) always leads to infinitely many resonances. Indeed, if, say,
$\kappa_0=0$, the resonance condition \eqref{eq:rescond} with any $\alpha$ can be
satisfied by taking $\beta_1=\alpha$ and $\beta_2=\beta_3=\dots=\beta_n=0$.

To illustrate this point let us consider an example with only two coupling
constants $y^0$, $y^1$, such that $\kappa_0=0$ and $\kappa_1=2$. I will use
suggestive notations $y^0=\lambda$ and $y^1=M^2$. Then in the flow equation for
$\lambda$ all terms $\lambda^n$ are resonant, so the coordinates exist in which
this equation takes the form
\begin{equation}\label{eq:normallambda}
\frac d{dl}\,\lambda=\sum_{n=2}^\infty b_n\,\lambda^n=-\beta(\lambda)\,.
\end{equation}
Similarly, in the equation for $M^2$ all terms of the form $M^2\lambda^n$ are
resonant, and in canonical coordinates one can write
\begin{equation}\label{eq:normalM}
\frac d{dl}\,M^2=2M^2-\sum_{n=1}^\infty\gamma_n\,M^2\lambda^n=\bigl(2-\gamma_{\varphi^2}(\lambda)\bigr)\,M^2\,.
\end{equation}
Note that with these notations these equations have exactly the form of the
coupling and the mass ``flow'' equations associated with the Callan-Symanzik
equation of the $d=4$ $\varphi^4$ theory.

\begin{remark}
In fact, in this important example a freedom still remains for further changes of
coordinates $\lambda\to\tilde\lambda(\lambda)$, $M^2\to\tilde M^2=M^2f(\lambda)$,
which doesn't change the above ``canonical'' form of the flow equations.
\end{remark}

\begin{exercise}
Check (to the order $n=$ when you get bored, or see the pattern) that all terms
$\sim\lambda^n$ with $n\ge4$ can be eliminated order by order, by suitable choice
of coefficients $a_n$ in
\[
\tilde\lambda=\lambda+a_3\,\lambda^3+a_4\,\lambda^4+\cdots\,.
\]
Further, check that suitable redefinition of $M^2\to\tilde M^2$,
\[
\tilde M^2=M^2\,\bigl(1+c_1\,\lambda+c_2\,\lambda^2+\cdots\bigr)
\]
allows one to eliminate the terms with $n\ge2$ in the r.h.s.\ of
\eqref{eq:normalM}. In reference to renormalization of the $\varphi^4$ theory at
$d=4$, these results imply that a scheme exists in which all $n$th order
corrections to the beta-function vanish starting from $n=4$, and all corrections
to $\gamma_{\varphi^2}$ vanish starting from $n=2$. At the same time the
coefficients $b_2$ and $b_3$, and $\gamma_1$, are scheme-independent.
\end{exercise}

Note that the statement of the Poincar\'e-Dulac theorem concerns formal power
series. Existence of ``canonical coordinates'', even locally, depends on
convergence of the series which in general is not guaranteed. For general RG flow
in the full space $\Sigma$ the situation is indeed difficult. One of the problems
lies again in the ``resonance denominators''
$\kappa_\alpha-\kappa_{\beta_1}-\dots-\kappa_{\beta_n}$ which the coefficients $a$
in the series \eqref{eq:normaltr} have. When both positive and negative
$\kappa_\alpha$ are present, even if no resonance conditions \eqref{eq:rescond} are
satisfied exactly, they often are ``almost satisfied'', in which case some
denominators come very close to zero giving rise to anomalously large coefficients
in \eqref{eq:normaltr}.

The situation becomes much more favorable if, instead of the flow in the full
space $\Sigma$, we concentrate attention on the RG flow in an unstable manifold
$U(\mathcal A_*)$ associated with some fixed point $\mathcal A_*$. First,
$U(\mathcal A_*)$ typically has finite dimensionality (as our study of
$\mathcal A_{*,G}$ illustrates). Also, by the definition of unstable manifold all
corresponding $\kappa$ are either positive (relevant) or zero (marginal). In this
case the resonance denominators are under much better control and the series
\eqref{eq:normaltr} defining the canonical coordinates usually converges in finite
neighborhood of $\mathcal A_*$. Recall that $U(\mathcal A_*)$ is a part of
$\Sigma(\infty)$, i.e.\ it describes finite local QFT. For local QFT generated by
the fixed point $\mathcal A_*$ the canonical coordinates do exist.

\subsection*{Callan-Symanzik equation revisited}

The preceding analysis allows one to re-derive the Callan-Symanzik equation which
we have obtained before on the basis of renormalized perturbation theory.

Let $\mathcal A_*$ be a fixed point. Assuming that its unstable manifold is
$n$-dimensional, let $\{\lambda^i\}$, $i=1,\dots,n$ be coordinates on
$U(\mathcal A_*)$, such that $\mathcal A_*$ corresponds to $\lambda=0$. This
unstable manifold describes $n$-parameter family of finite local field theories. I
will denote the corresponding correlation functions as
\[
\bigl\langle\varphi(x_1)\cdots\varphi(x_N)\bigr\rangle_\lambda\,,
\]
indicating their dependence on $n$ parameters $\lambda^i$.

In the coordinates $\lambda$ the RG flow equation can be written as
\begin{equation}\label{eq:flowlambdai}
\frac d{dl}\,\lambda^i=-\beta^i(\lambda)\,,
\end{equation}
where $-\beta^i$ are components of the vector field $B^\alpha$ on
$U(\mathcal A_*)$ (by definition of $U$, $B(U)$ is tangent to $U$), the minus sign
is taken to match our previous notation for the beta-function. According to the
definition of the Wilson's RG, if $\lambda(l)$ is a solution of
\eqref{eq:flowlambdai} with some initial condition $\lambda(0)=\lambda$, the field
theories corresponding to $\lambda$ and $\lambda(l)$ are related by the scale
transformation $x\to x/L$. In particular, one can write for the correlation
functions
\[
\bigl\langle\varphi(x_1)\cdots\varphi(x_N)\bigr\rangle_{\lambda(l)}=Z^{\frac N2}(L)\,\bigl\langle\varphi(Lx_1)\cdots\varphi(Lx_N)\bigr\rangle_{\lambda(0)}\,,
\]
where $L=e^l$. Taking infinitesimal transformation $l=\delta l$, $L=1+\delta l$,
$\lambda(l)=\lambda-\beta(\lambda)\,\delta l$, we have
\begin{equation}\label{eq:CSgeneral}
\Bigl[x\frac\partial{\partial x}+N\,D(\lambda)+\sum_{i=1}^n\beta^i(\lambda)\,\frac\partial{\partial\lambda^i}\Bigr]\bigl\langle\varphi(x_1)\cdots\varphi(x_N)\bigr\rangle_\lambda=0\,,
\end{equation}
where $x\frac\partial{\partial x}$ is the shorthand notation for
\[
x\frac\partial{\partial x}=\sum_{k=1}^Nx_k^\mu\,\frac\partial{\partial x_k^\mu}\,,
\]
and
\[
D(\lambda)=L\frac d{dL}\,Z^{\frac12}(L)\Big|_{L=1}\,.
\]
Note that $D(0)=D_\varphi$, the dimension of $\varphi$ at the fixed point
$\mathcal A_*$. $D(\lambda)$ can be written as
\[
D(\lambda)=D_\varphi+\gamma(\lambda)\,,
\]
so that $\gamma(0)=0$.

The equation \eqref{eq:CSgeneral} generalizes straightforwardly to correlation
functions involving composite fields $\mathcal O_\alpha$. Repeating the same
analysis we obtain
\begin{multline}\label{eq:CScomposite}
\sum_{k=1}^N\bigl\langle\mathcal O_{\alpha_1}(x_1)\cdots(\mathcal D\mathcal O)_{\alpha_k}(x_k)\cdots\mathcal O_{\alpha_N}(x_N)\bigr\rangle_\lambda\\
+\sum_{i=1}^n\beta^i(\lambda)\,\frac\partial{\partial\lambda^i}\bigl\langle\mathcal O_{\alpha_1}(x_1)\cdots\mathcal O_{\alpha_N}(x_N)\bigr\rangle_\lambda=0\,,
\end{multline}
where
\[
(\mathcal D\mathcal O)_\alpha(x)=x^\mu\frac\partial{\partial x^\mu}\,\mathcal O_\alpha(x)+D_\alpha^\beta(\lambda)\,\mathcal O_\beta(x)\,,
\qquad\text{and}\qquad
D_\alpha^\beta(\lambda)=-L\frac d{dL}\,z_\alpha^\beta(L)\Big|_{L=1}\,;
\]
$z_\alpha^\beta(L)$ describes the RG transformation of the composite fields
$\mathcal O_\alpha$ (see Sect.~18). Again,
\[
D_\alpha^\beta=D_\alpha^\beta(\lambda)\big|_{\lambda=0}
\]
coincides with the matrix of dimensions at $\mathcal A_*$, whose eigenvalues are
the dimensions of the fields $\Phi_\alpha$.

We see that the CS equation is very general consequence of Wilson's RG theory and
not just the property of the perturbation theory around the Gaussian fixed point
(i.e.\ free field theory).
